\chapter{Object (swallowed accidentally)}

\section*{Definition and Overview}
Accidental swallowing of a foreign object, clinically termed foreign body ingestion (FBI), refers to the entry of non-food items, large or poorly chewed food boluses, or dental prostheses into the gastrointestinal tract. While the majority of ingested objects pass through the digestive system without incident, certain items pose significant clinical hazards due to their size, shape, chemical composition, or electrical properties. Lodgement can occur at sites of anatomical narrowing, leading to mucosal injury, obstruction, or perforation. Foreign body ingestion is a critical medical concern requiring prompt triaging to distinguish between benign, self-limiting cases and life-threatening emergencies.

\section*{Epidemiology}
Foreign body ingestion occurs most frequently in the paediatric population, with peak incidence between six months and three years of age, representing approximately 80\% of all clinical cases. In children, coins, small toys, jewellery, and button batteries are the most commonly ingested items. In adults, ingestion is typically accidental and frequently involves fish or poultry bones, food boluses (often associated with underlying oesophageal pathology), or dental components such as crowns and partial dentures. High-risk adult populations include individuals with intellectual disability, psychiatric disorders, cognitive impairment, or those involved in illicit drug packaging (body packing). 

\section*{Aetiology and Pathophysiology}
The physiological mechanisms and risk factors contributing to foreign body ingestion and subsequent tissue damage include:
\begin{itemize}
    \item \textbf{Paediatric developmental behaviours:} Young children explore their environment through mouthing behaviours, placing small household items at risk of accidental ingestion.
    \item \textbf{Underlying oesophageal pathology:} In adults, food bolus impactions are frequently associated with anatomical or functional abnormalities, such as eosinophilic oesophagitis, peptic strictures, Schatzki rings, or motility disorders.
    \item \textbf{Anatomical narrowing:} Lodgement typically occurs at zones of physiological constriction, including the cricopharyngeal sphincter (upper oesophageal sphincter), the level of the aortic arch, the left main bronchus crossing, and the lower oesophageal sphincter.
    \item \textbf{Liquefactive necrosis from button batteries:} Lodged button batteries generate an electrical current that hydrolyses tissue fluids, creating hydroxide ions. This leads to rapid alkaline liquefactive necrosis of the oesophageal mucosa, which can penetrate deep tissues within two hours.
    \item \textbf{Mechanical trauma:} Sharp objects, such as pins, needles, and bones, can directly puncture the mucosal lining, causing focal inflammation, abscess formation, or free perforation into the mediastinum or peritoneal cavity.
    \item \textbf{Magnetic attraction:} Ingestion of multiple magnets can lead to their attraction across different loops of the bowel, compressing the intervening intestinal wall and causing pressure necrosis, volvulus, fistula formation, or perforation.
\end{itemize}

\section*{Clinical Features}
The clinical presentation of foreign body ingestion varies according to the child's or adult's age, the nature of the object, and its anatomical location:
\begin{itemize}
    \item Dysphagia (difficulty swallowing) or aphagia (complete inability to swallow)
    \item Odynophagia (painful swallowing) or retrosternal chest pain
    \item Drooling, pooling of saliva, or inability to clear secretions
    \item Choking, coughing, wheezing, stridor, or acute respiratory distress (indicating airway compression or aspiration)
    \item Foreign body sensation in the throat or chest
    \item Abdominal pain, abdominal distension, vomiting, or haematemesis (typically indicating gastric or intestinal complications)
    \item Irritability, refusal to feed, or unexplained fever in non-verbal children
\end{itemize}

\section*{Differential Diagnosis}
When evaluating a patient with suspected foreign body ingestion, the following clinical entities must be considered:
\begin{itemize}
    \item \textbf{Aspiration of a foreign body:} Displacement of an object into the trachea or mainstem bronchi rather than the oesophagus, requiring urgent bronchoscopy.
    \item \textbf{Oesophageal motility disorders:} Conditions like achalasia or distal oesophageal spasm, which present with acute dysphagia and chest pain.
    \item \textbf{Gastro-oesophageal reflux disease (GORD):} Reflux-induced oesophagitis or globus pharyngeus mimicking a persistent foreign body sensation.
    \item \textbf{Acute epiglottitis or croup:} Inflammatory airway disorders in children presenting with drooling, stridor, and dysphagia.
    \item \textbf{Peptic ulcer disease or gastritis:} Causes of acute epigastric pain and haematemesis without a history of ingestion.
\end{itemize}

\section*{When to Seek Medical Care}
Medical assessment should be sought immediately for any individual who has swallowed a high-risk object (such as a button battery, magnet, or sharp item) or who is exhibiting symptoms of gastrointestinal or respiratory distress.

\textbf{Call 000 (or local emergency services) for:}
\begin{itemize}
    \item Any suspected or confirmed ingestion of a button battery or multiple magnets
    \item Acute breathing difficulties, stridor, choking, or persistent coughing
    \item Complete inability to swallow saliva, leading to constant drooling
    \item Severe, localized pain in the neck, chest, or abdomen
    \item Vomiting blood or material resembling coffee grounds, or passing black, tarry stools
\end{itemize}

\section*{Diagnosis and Investigations}
A structured diagnostic approach is necessary to locate the object and identify potential complications:
\begin{itemize}
    \item \textbf{Plain radiography:} Biplanar X-rays (anteroposterior and lateral views) of the neck, chest, and abdomen are the initial imaging standard. Radiopaque objects (e.g., coins, metal keys) are easily identified, and button batteries display a characteristic double-contour ``halo'' sign on AP view and a step-off on lateral view.
    \item \textbf{Computed tomography (CT) scan:} Indicated for radiolucent objects (e.g., plastic parts, wood, fish bones), suspected perforation, or complex cases where clinical suspicion remains high despite negative X-rays.
    \item \textbf{Diagnostic upper endoscopy:} The gold standard for directly visualising the mucosal surfaces and locating objects in the oesophagus, stomach, or proximal duodenum.
    \item \textbf{Handheld metal detector:} A useful, non-invasive screening tool in paediatric settings to trace the progress of confirmed metallic, non-hazardous foreign bodies like coins.
\end{itemize}

\section*{Management}
The management strategy depends on the patient's symptoms, the anatomical location of the foreign body, and the risk profile of the object:
\begin{itemize}
    \item \textbf{Urgent endoscopic retrieval:} Indicated for any foreign body in the oesophagus causing complete obstruction, button batteries lodged in the oesophagus (within two hours), sharp objects, or multiple magnets within reach of an upper endoscope.
    \item \textbf{Conservative monitoring:} Appropriate for asymptomatic patients who have ingested small, blunt, non-hazardous objects (such as coins) that have already passed into the stomach. Stools should be monitored, and follow-up radiographs may be performed if the object has not passed within one to two weeks.
    \item \textbf{Surgical intervention (laparoscopy or laparotomy):} Required if the object fails to progress beyond the duodenum after several days, if the patient develops signs of bowel obstruction, or if there are clinical features of peritonitis or perforation.
    \item \textbf{Avoidance of emetics:} Emetics (drugs that induce vomiting) should not be administered, as refluxing the object can cause airway aspiration or secondary mucosal lacerations.
\end{itemize}

\section*{Complications}
Failure to promptly diagnose or treat hazardous ingestions can lead to severe outcomes:
\begin{itemize}
    \item Oesophageal perforation and secondary mediastinitis
    \item Aortoesophageal fistula, causing catastrophic, life-threatening haemorrhage
    \item Tracheoesophageal fistula, resulting in recurrent aspiration and lung infections
    \item Small or large bowel obstruction
    \item Intestinal perforation, leading to peritonitis and sepsis
    \item Mucosal ulceration, stricture formation, and long-term dysphagia
\end{itemize}

\section*{Prognosis and Outcomes}
The prognosis is excellent for the vast majority of patients. Approximately 80\% to 90\% of ingested foreign bodies pass spontaneously through the gastrointestinal tract within a few days to a week without requiring intervention. Endoscopic removal is successful in 10\% to 20\% of cases, while surgical intervention is required in less than 1\% of patients. However, if a button battery remains lodged in the oesophagus, tissue damage can progress rapidly, and delayed intervention is associated with significant morbidity and potential mortality.

\section*{Prevention and Self-Care}
\begin{itemize}
    \item Store all button batteries, coins, pins, and small household items in secure, childproof containers out of reach of children.
    \item Ensure battery compartments on toys, remote controls, and keys are screwed shut or taped securely.
    \item Cut food into small, manageable pieces, chew thoroughly before swallowing, and avoid talking or laughing while eating.
    \item Avoid holding pins, nails, or screws between the lips or teeth while performing domestic or mechanical tasks.
    \item Keep dentures well-fitted and seek dental review if they become loose or damaged.
\end{itemize}

\section*{Sources and Further Reading}
The following reputable organisations provide further information and clinical guidelines on foreign body ingestion:
\begin{itemize}
    \item Healthdirect Australia. \textit{Accidental swallowing of foreign objects.} https://www.healthdirect.gov.au/
    \item Royal Children's Hospital Melbourne. \textit{Clinical Practice Guidelines: Foreign body ingestion.} https://www.rch.org.au/
    \item Mayo Clinic. \textit{Foreign object swallowed: First aid.} https://www.mayoclinic.org/
    \item NHS. \textit{What to do if someone swallows a foreign object.} https://www.nhs.uk/
    \item MSD Manual (Consumer Version). \textit{Foreign Bodies in the Gastrointestinal Tract.} https://www.msdmanuals.com/
\end{itemize}
